\section{有序集}

\begin{enumerate}
\def\labelenumi{\arabic{enumi}.}
\tightlist
\item
  一个关系 \(\omega\{x, y\}\) 在两个一般元素之间 \(x, y\) 集合 E 上的一个关系被称为 E 上的序关系，如果它满足以下两个条件：
\end{enumerate}

\begin{enumerate}
\def\labelenumi{(\alph{enumi})}
\item
  关系“ \(\omega \{x, y\}\) 并且 \(\omega \{y, z\}\) '' 蕴含 \(\omega \{x, z\}\) （传递性）。
\item
  关系“ω\{x, y\} 且 ω\{y, x\}'' 等价于 “x = y''。条件 (b) 蕴含 ω 是自反的。
\end{enumerate}

设 C 是 \(E \times E\) 的子集，由关系 \(\omega\{x, y\}\) 作为对 \((x, y)\) 的性质来定义。则条件 (a) 和 (b) 分别等价于以下内容： \((a') \quad C \circ C \subset C\) ; \((b') \quad C \cap \overline{C} = \Delta\) 。这些蕴含 \(C \circ C = C\) .

当我们考虑集合 E 上的一个特定序关系时，我们说 E 由该关系排序，并且该关系定义了一个序结构（参见 §8）（或一个排序）在 E 上。

如果 \(\omega\{x,y\}\) 是 E 上的序关系，则 \(\omega\{y,x\}\) 也是。这两个序关系，以及它们定义的排序，被称为互为相反关系。

令 \(\omega\{x, y\}\) 是 E 上的一个序关系，设 A 是 E 的任意子集。则关系 \(\omega\{x, y\}\) （A 的两个一般元素 \(x, y\) 之间的）是 A 上的一个序关系，且它在A上定义的序称为由 \(\omega\{x, y\}\) 在E上定义的序所诱导的。E上给定的序称为A上诱导序的一个扩张。

一个自反且传递的关系 \(\varpi\{x, y\}\) （E 的两个一般元素 \(x, y\) 之间的）称为E上的一个预序关系。关系“ \(\varpi\{x, y\}\) 并且 \(\varpi\{y, x\}\) '' 是 E 上的一个等价关系 R，并且 \(\varpi\{x, y\}\) 与此关系（在 \(x\) 和 \(y\) 中）相容。通过过渡到商（关于 \(x\) 和 \(y\) ）， \(\varpi\{x, y\}\) 在集合 E/R 上诱导出一个序关系，称为与 \(\varpi\{x, y\}\) 相关联的序关系。一个配备了预序关系的集合称为预序集。

\begin{enumerate}
\def\labelenumi{\arabic{enumi}.}
\setcounter{enumi}{1}
\tightlist
\item
  包含关系“Y ⊂ X”是子集集合上的一个序关系。 \(\mathfrak{P}(\mathbf{E})\) 对于任意集合E。
\end{enumerate}

若E和F是两个集合，它们可能相同也可能不同，则“g延拓f”这一关系是E的子集到F的所有映射之集合上的一个序关系。

自然整数集合 N (*) 由关系 `` \(x \leqslant y\) '' 排序。3. 与最后一个例子类似，当集合 E 由关系 \(\omega\{x, y\}\) 排序时，通常方便将关系记为 \(\omega\{x, y\}\) 由 \(x \leqslant y\) ，或 \(y \geqslant x\) ；这些关系读作 `` \(x\) 小于 \(y\) ''，或`` \(y\) 大于 \(x\) ''（†）。关系`` \(x < y\) ”和“ \(y > x\) ''（读作`` \(x\) 严格小于 \(y\) ''，或`` \(y\) 严格大于 \(x\) ''）按定义等价于`` \(x \leqslant y\) 并且 \(x \neq y\) ''.

关系``x ≤ y''等价于``x \textless{} y 或 x = y''。关系``x ≤ y 且 y \textless{} z''蕴含``x \textless{} z''；类似地，``x \textless{} y 且 y ≤ z''蕴含``x \textless{} z''。

\begin{enumerate}
\def\labelenumi{\arabic{enumi}.}
\setcounter{enumi}{3}
\tightlist
\item
  集合 E 的子集 X，由关系`` \(x \leqslant y\) ''排序，如果对于所有 \(x \in X\) 以及所有 \(y \in X\) ，我们有要么 \(x \leqslant y\) 要么 \(y \leqslant x\) （或等价地，要么 x \textless{} y 或 x = y 或 x \textgreater{} y，这三个关系互斥），则称该子集被此关系全序化。
\end{enumerate}

有序集合的空子集总是全序的。可能发生整个集合 E 是全序的情况，例如集合 N 和关系 \(x \leqslant y\) .

全序集的每个子集在诱导序下都是全序的。

在有序集合 E 中，如果 a 和 b 是两个元素使得 \(a \leqslant b\) ，则 E 中由满足 \(a \leqslant x \leqslant b\) 的元素组成的子集称为以 a 为左端点、b 为右端点的闭区间，并记为 \([a, b]\) 。所有满足 \(x \in E\) 且 \(a \leqslant x < b\) （相应地 \(a < x \leqslant b\) ）的集合称为以 a 和 b 为端点的右半开（相应地，左半开）区间，并记为 \([a, b[\) （相应地 {]}a, b{]}）。所有满足 \(x \in E\) （且 a \textless{} x \textless{} b）的 x 组成的集合称为以 a 和 b 为端点的开区间，并记为 {]}a, b{]} 。注意，闭区间永远非空.

所有满足 \(x \in \mathbf{E}\) 且 \(x \leqslant a\) （相应地 \(x < a\) ）的 x 组成的集合称为以 \(a\) 为右端点的左无界闭（相应地，开）区间，并记作 \(] \leftarrow, a]\) （相应地 \(] \leftarrow, a[\) ）；同样地，所有满足 \(x \in \mathbf{E}\) 且 \(x \geqslant a\) （相应地 \(x > a\) ）的元素 x 组成的子集称为以 \(a\) 为左端点的右无界闭（相应地，开）区间，并记作 \([a, \rightarrow[\) （相应地 \(]a, \rightarrow[]\) ）。最后，我们将 E 本身视为双向无界的开区间，并如此将其记为 \(] \leftarrow, \rightarrow]\) .

\begin{enumerate}
\def\labelenumi{\arabic{enumi}.}
\setcounter{enumi}{4}
\tightlist
\item
  如果 X 是有序集合 E 的子集，则至多存在一个 X 中的元素 \(a\) 使得 \(a \leqslant x\) 对于所有 \(x \in \mathbf{X}\) ；当存在具有此性质的元素 \(a\) 时，它被称为 X 的最小元素。同样地，至多存在一个元素 \(b\) 使得 \(x \leqslant b\) 对于所有 \(x \in \mathbf{X}\) ; \(b\) ，如果存在，则被称为 X 的最大元素。
\end{enumerate}

在全序集合 E 中，每个有限非空子集都有一个最大元素和一个最小元素，分别称为该子集的最大元和最小元。

一个有序集合，如果其每个非空子集都有最小元素，则称为良序的。自然整数集合 N 是良序的；N 的子集有最大元素当且仅当它是有限且非空的。利用选择公理可以证明，在每个集合上都存在一个良序（策梅洛定理）。

在集合 \(\mathfrak{P}(\mathbf{E})\) （集合 \(\mathbf{E}\) 的子集所成的，按包含关系排序）中，一个子集 \(\mathfrak{F}\) （ \(\mathfrak{P}(\mathbf{E})\) 的）有最小元素当且仅当 \(\mathfrak{F}\) 中这些集合的交集属于 \(\mathfrak{F}\) ，并且该交集就是最小元素。类似地， \(\mathfrak{F}\) 有最大元素当且仅当 \(\mathfrak{F}\) 中这些集合的并集属于 \(\mathfrak{F}\) ，并且该并集就是最大元素。

有序集 E 的子集 X 被称为在 E 中共尾（相应地，共首），如果对于每个 \(y \in E\) , 存在 \(x \in X\) 使得 \(y \leqslant x\) （相应地 \(y \geqslant x\) ）。说有序集 E 有最大（相应地，最小）元素，意味着存在 E 的一个由单个元素组成的共尾（相应地，共首）子集。

\begin{enumerate}
\def\labelenumi{\arabic{enumi}.}
\setcounter{enumi}{5}
\item
  设 X 是有序集 E 的一个子集。每个 \(x \in X\) 使得不存在元素 \(z \in X\) 满足 z \textless{} x 的元素称为 X 的极小元素。每个 \(y \in X\) 使得不存在元素 \(z \in X\) 满足 z \textgreater{} y 的元素称为 X 的极大元素。极大元素（或极小元素）的集合可能为空；也可能无限。如果 X 有最小元素 a，则 a 是 X 唯一的极小元素；同样，如果 X 有最大元素 b，则 b 是 X 唯一的极大元素。
\item
  设 X 是有序集 E 的一个子集。如果元素 \(x \in E\) 是这样的 \(z \leqslant x\) 对于所有 \(z \in X\) ，则 z 称为 X 的上界。类似地，元素 \(y \in E\) 使得 \(z \geqslant y\) 对于所有 \(z \in X\) 称为 X 的下界。
\end{enumerate}

子集 X 的上界集合（或下界集合）可能为空。上界（相应地，下界）集合非空的子集 X 被称为上有界（相应地，下有界）。既有上界又有下界的集合简称为有界。大于 X 的上界的每个元素都是 X 的上界；小于 X 的下界的每个元素都是 X 的下界。

如果子集 X 的上界集合有最小元素 a，则 a 称为 X 的最小上界或上确界。如果 X 的下界集合有最大元素 b，则 b 称为 X 的最大下界或下确界。如果这些界存在，则由定义它们是唯一的，并且分别记为 \(\sup_{E} X\) （或 \(\sup X\) ), \(\inf_{E} X\) （或 \(\inf X\) ). 若X有最大元，则它必为X的上确界；若X有最小元，则它必为X的下确界。反之，若X的上确界（相应地，下确界）存在且属于X，则它必为X的最大元（相应地，最小元）。

让 \(f\) 是从集合 \(A\) 到 \(E\) 的一个映射。如果 \(f(A)\) 在 \(E\) 中有上界（相应地，有下界，有界），那么 \(f\) 被称为上有界（相应地，下有界，有界）。如果 \(f(A)\) 在 \(E\) 中有最小上界（相应地，有最大下界），这个界被称为 \(f\) 的最小上界（相应地，最大下界）并且用 \(\sup_{x \in A} f(x)\) （相应地 \(\inf_{x \in A} f(x)\) ).

\begin{enumerate}
\def\labelenumi{\arabic{enumi}.}
\setcounter{enumi}{7}
\tightlist
\item
  一个预序集 E，如果 E 的每个有限非空子集都有上界（分别地，下界），则称其为右定向的（或定向的）（分别地，左定向的）。
\end{enumerate}

一个有序集 E，如果 E 的每个有限非空子集都有最小上界和最大下界，则称其为格。

任意集合的子集族，按包含关系排序，构成一个格。每个全序集都是格。

\begin{enumerate}
\def\labelenumi{\arabic{enumi}.}
\setcounter{enumi}{8}
\tightlist
\item
  一个有序集 E 被称为归纳的，如果它满足以下条件：E 的每个全序子集都有上界。
\end{enumerate}

集合 \(\mathfrak{P}(\mathbf{E})\) ，按包含关系排序，是归纳的。将集合 E 的子集映射到集合 F 的映射所构成的集合，在按 f 与 g 之间的关系“g 延拓 f”排序时，也是归纳的。

归纳集的任意子集一般不一定是归纳的。但如果 \(a\) 是归纳集的任意元素 \(\mathbf{E}\) ，则子集 \(\mathbf{E}\) 由所有元素组成 \(x \in \mathbf{E}\) 使得 \(x \geqslant a\) 是归纳的。

\begin{enumerate}
\def\labelenumi{\arabic{enumi}.}
\setcounter{enumi}{9}
\tightlist
\item
  以下命题借助选择公理证明，被称为佐恩引理：
\end{enumerate}

每个归纳有序集至少有一个极大元。

\begin{enumerate}
\def\labelenumi{\arabic{enumi}.}
\setcounter{enumi}{10}
\tightlist
\item
  在子集的集合中 \(\mathfrak{P}(\mathbf{E})\) 的集合 \(\mathbf{E}\) ，按包含关系排序，一个集合的最小上界 \(\mathfrak{F}\) 子集的 \(\mathbf{E}\) 是以下集合的并集 \(\mathfrak{F}\) 应用佐恩引理可得到以下结果：
\end{enumerate}

如果 \(\mathfrak{F}\) 是集合 E 的子集族，使得对于每个子集 \(\mathfrak{G}\) 的 \(\mathfrak{F}\) 它由包含关系全序化，则这些集合的并集 \(\mathfrak{G}\) 属于 \(\mathfrak{F}\) ，那么 \(\mathfrak{F}\) 至少有一个极大元素（即E的一个子集，它属于 \(\mathfrak{F}\) 但不包含于任何其他属于 \(\mathfrak{F}\) ).

一个由集合 E 的子集组成的集合 \(\mathfrak{F}\) 被称为具有有限特征，如果性质“ \(\mathbf{X} \in \mathfrak{F}\) ”等价于性质“X 的每个有限子集都属于 \(\mathfrak{F}\) ”。有了这个定义，我们有以下定理：

E 的每个具有有限特征的子集族至少有一个极大元素。

\begin{enumerate}
\def\labelenumi{\arabic{enumi}.}
\setcounter{enumi}{11}
\tightlist
\item
  一个映射 \(f\) （从子集 \(\mathbf{A}\) （属于有序集合 \(\mathbf{E}\) ）到有序集 \(\mathbf{F}\) ）被称为递增（相应地，递减）的，如果关系 \(x \leqslant y\) （在 \(\mathbf{A}\) 的典型元素之间成立）蕴含 \(f(x) \leqslant f(y)\) （相应地 \(f(y) \leqslant f(x)\) ). \(\mathbf{A}\) 上的每个常值函数因此既递增又递减，且若 \(\mathbf{A}\) 是（右向或左向）有向的，则逆命题成立.
\end{enumerate}

一个映射 \(f\) 被称为严格递增（相应地，严格递减）的，如果关系 \(x < y\) 蕴含 \(f(x) < f(y)\) （相应地 \(f(x) < f(y)\) ).

若E是全序的，则E的子集A到有序集F的任意严格递增（或严格递减）映射都是单射。

如果 I 是一个有序指标集，则集合 E 的子集的一个族 \((\mathbf{X}_{\iota})_{\iota\in\mathbf{I}}\) 被称为递增的（相应地，递减的），如果映射 \(\iota\to X_{\iota}\) （从 I 到 \(\mathfrak{P}(\mathbf{E})\) ，后者按包含关系排序）是递增的（相应地，递减的）。

\begin{enumerate}
\def\labelenumi{\arabic{enumi}.}
\setcounter{enumi}{12}
\tightlist
\item
  设 I 是一个有向集，并令 \((\mathrm{E}_{\alpha})_{\alpha\in\mathbf{I}}\) 是由I指标化的一族集合。对于每一对 \((\alpha,\beta)\) （指标集合I中的，满足 \(\alpha\leqslant\beta\) ），令 \(f_{\beta\alpha}\) 为从 \(E_{\alpha}\) 到 \(E_{\beta}\) 的映射，并假设关系 \(\alpha\leqslant\beta\leqslant\gamma\) 蕴含 \(f_{\gamma\alpha}=f_{\gamma\beta}\circ f_{\beta\alpha}\) .
\end{enumerate}

让 \(\mathbf{F}\) 是集合族 \((\mathbf{E}_{\alpha})_{\alpha \in \mathbf{I}}\) 的和。为了语言上的简便，我们将 \(\mathbf{E}_{\alpha}\) 与 \(\mathbf{F}\) 中相应的子集等同起来。给定两个元素 \(x\) 和 \(y\) （属于 \(\mathbf{F}\) ），让 \(\mathbb{R}\{x, y\}\) 为如下关系（其中 \(\alpha\) 和 \(\beta\) 表示 I 中满足 \(x \in \mathbf{E}_{\alpha}\) 并且 \(y \in \mathbf{E}_{\beta})\) :

“存在 \(\gamma \in I\) 使得 \(\gamma \geqslant \alpha\) 并且 \(\gamma \geqslant \beta\) 并且 \(f_{\gamma \alpha}(x) = f_{\gamma \beta}(y) "\) .

则 R 是 F 上的一个等价关系。设 E 为商集 F/R，并设 f 为典范映射 \(F \rightarrow F/R\) 。集合 E 称为族 \((\mathbf{E}_{\alpha})_{\alpha \in \mathbf{I}}\) 关于映射族 \((f_{\beta\alpha})\) 的直接极限，且限制 \(f_{\alpha}\) （即 f 到 \(E_{\alpha}\) 上的限制）称为 \(E_{\alpha}\) 到 E 中的典范映射。我们有 \(f_{\beta} \circ f_{\beta\alpha} = f_{\alpha}\) 每当 \(\alpha \leqslant \beta\) 。我们记

\[
\mathbf {E} = \lim _ {\rightarrow} (\mathrm{E} _ {\alpha}, f _ {\beta \alpha})
\]

或简言之 \(E = \lim_{\alpha} E_{\alpha}\) 当无混淆风险时。滥用语言地，称对 \(((E_{\alpha}), (f_{\beta\alpha}))\) 为相对于 I 的集合直接系统。

语言，称对 \(((\mathbf{E}_{\alpha}), (f_{\beta \alpha}))\) 为相对于 I 的集合直接系统。若 \(f_{\beta \alpha}\) 是单射的，则 \(f_{\alpha}\) 是单射的。在这种情况下，我们通常将 \(\mathbf{E}_{\alpha}\) 与 \(f_{\alpha}(\mathbf{E}_{\alpha})\) 等同起来，因此将 E 视为 \(\mathbf{E}_{\alpha}\) 的并集。反之，若一个集合 \(\mathbf{E}'\) 是子集的某个族 \((\mathbf{E}_{\alpha}')_{\alpha \in \mathbf{I}}\) 的并集，使得关系 \(\alpha \leqslant \beta\) 蕴含 \(\mathbf{E}_{\alpha}' \subset \mathbf{E}_{\beta}'\) ，并且如果（对于 \(\alpha \leqslant \beta\) ） \(j_{\beta \alpha}\) 表示 \(\mathbf{E}_{\alpha}'\) 到 \(\mathbf{E}_{\beta}'\) 的典范单射，那么我们可以将 \(\lim (\mathbf{E}_{\alpha}', j_{\alpha \beta})\) 与 \(\mathbf{E}'\) 等同起来，并且将 \(\mathbf{E}_{\alpha}'\) 到 \(\lim (\overrightarrow{\mathbf{E}_{\alpha}', j_{\beta \alpha}})\) 中的典范映射与 \(\mathbf{E}_{\alpha}'\) 到 E 中的典范单射等同起来。

更一般地，设 \((\mathbf{E}_{\alpha}, f_{\beta \alpha})\) 是相对于 I 的集合直接系统，并且对每个 \(\alpha \in I\) 让 \(g_{\alpha}\) 是从 \(\mathbf{E}_{\alpha}\) 到一个集合 \(\mathbf{E}'\) 的映射，使得关系 \(\alpha \leqslant \beta\) 蕴含 \(g_{\beta} \circ f_{\beta \alpha} = g_{\alpha}\) 。则存在唯一的映射 \(g\) （从 \(\mathbf{E} = \lim_{\rightarrow} \mathbf{E}_{\alpha}\) 到 \(\mathbf{E}'\) ）使得 \(g_{\alpha} = g \circ f_{\alpha}\) 对于所有 \(\alpha \in I\) . 映射 \(g\) 是满射当且仅当 \(\mathbf{E}'\) 是 \(g_{\alpha}(\mathbf{E}_{\alpha})\) 的并集. 映射 \(g\) 是单射当且仅当，对每个 \(\alpha \in I\) ，关系 \(x \in \mathbf{E}_{\alpha}\) , \(y \in \mathbf{E}_{\alpha}, g_{\alpha}(x) = g_{\alpha}(y)\) 意味着存在 \(\beta \geqslant \alpha\) 使得 \(f_{\beta \alpha}(x) = f_{\beta \alpha}(y)\) 。如果 \(g\) 是双射， \(\mathbf{E}'\) 有时与 \(\mathbf{E}_{\alpha}\) 的直接极限等同起来.

让 \((\mathbf{A}_{\alpha}, \varphi_{\beta \alpha})\) 并且 \((\mathbf{B}_{\alpha}, \psi_{\beta \alpha})\) 是相对于同一指标集 I 的两个集合直接系统。设 \(\mathbf{A} = \varinjlim (\mathbf{A}_{\alpha}, \varphi_{\beta \alpha}), \mathbf{B} = \varinjlim (\mathbf{B}_{\alpha}, \psi_{\beta \alpha})\) ，并且对于每个 \(\alpha \in I\) 让 \(\varphi_{\alpha}\) （相应地 \(\psi_{\alpha}\) ）表示 \(\mathbf{A}_{\alpha}\) 到 A（相应地， \(\mathbf{B}_{\alpha}\) 到 B）的典范映射。对于每个 \(\alpha \in I\) 让 \(u_{\alpha}\) 是从 \(\mathbf{A}_{\alpha}\) 到 \(\mathbf{B}_{\alpha}\) 的映射，使得 \(u_{\beta} \circ \varphi_{\beta \alpha} = \psi_{\beta \alpha} \circ u_{\alpha}\) 每当 \(\alpha \leqslant \beta\) 。族 \((u_{\alpha})\) 称为 \((\mathbf{A}_{\alpha}, \varphi_{\beta \alpha})\) 到 \((\mathbf{B}_{\alpha}, \psi_{\beta \alpha})\) 的映射直接系统。在这些条件下，存在唯一的映射 \(u: \mathbf{A} \to \mathbf{B}\) 使得 \(u \circ \varphi_{\alpha} = \psi_{\alpha} \circ u_{\alpha}\) 对于所有 \(\alpha \in I\) . 映射 \(u\) 称为 \(u_{\alpha}\) 的直接极限，记作 \(u = \varprojlim u_{\alpha}\) ，前提是没有混淆的风险。设 \((C_{\alpha}, \theta_{\beta \alpha})\) 是相对于I的另一个集合直接系统，设 \((v_{\alpha})\) 是 \((B_{\alpha}, \psi_{\beta \alpha})\) 到 \((C_{\alpha}, \theta_{\beta \alpha})\) 的映射直接系统，并且让 \(v = \varprojlim v_{\alpha}\) 。则 \(\varinjlim (v_{\alpha} \circ u_{\alpha}) = v \circ u.\)

保持上述记号，设 \(D_{\alpha}=A_{\alpha}\times B_{\alpha}\) 并且 \(\omega_{\beta\alpha}=\varphi_{\beta\alpha}\times\psi_{\beta\alpha}\) 。那么族 \((\mathbf{D}_{\alpha},\omega_{\beta\alpha})\) 是集合的直接系统。设 \(\mathbf{D}=\lim(\mathbf{D}_{\alpha},\omega_{\beta\alpha})\) ，让 \(\omega_{\alpha}\) 是 \(\mathbf{D}_{\alpha}\) 到 \(\mathbf{D}\) 的典范映射，让 \(\mathbf{D}' = \mathbf{A} \times \mathbf{B}\) ，并且让 \(\omega_{\alpha}' = \varphi_{\alpha} \times \psi_{\alpha}\) 。则存在唯一的双射（称为典范的） \(f: \mathbf{D} \to \mathbf{D}'\) 使得 \(f \circ \omega_{\alpha} = \omega_{\alpha}'\) 对于所有 \(\alpha \in I\) . 我们通常将直接极限的乘积 \(\mathbf{D}'\) 与直接极限 \(\mathbf{D}\) （即诸乘积 \(\mathbf{D}_{\alpha}\) 的直接极限）等同起来.

设 J 是 I 的共尾子集；则 J 也是有向的。若 \(((\mathrm{E}_{\alpha})_{\alpha\in\mathbf{I}}, (f_{\beta\alpha})_{\alpha\in\mathbf{I},\beta\in\mathbf{I}})\) 是相对于I的集合直系统，则 \(((\mathrm{E}_{\alpha})_{\alpha\in\mathbf{J}}, (f_{\beta\alpha})_{\alpha\in\mathbf{J},\beta\in\mathbf{J}})\) 是关于J的集合的正向系统。设 \(E'\) 为其直极限，并令 \(f_{\alpha}'\) 是 \(E_{\alpha}\) 到 \(E'\) 的典范映射，对于所有 \(\alpha\in J\) 。则存在唯一的映射 \(g:E\to E'\) 使得 \(g(f_{\alpha}'(x))=f_{\alpha}(x)\) 对于所有 \(\alpha\in J\) 以及所有 \(x\in E_{\alpha}\) （其中 \(f_{\alpha}\) 表示 \(E_{\alpha}\) 到 E 中的典范映射）。此映射是双射，借助它 \(E'\) 通常与 E 等同。

\begin{enumerate}
\def\labelenumi{\arabic{enumi}.}
\setcounter{enumi}{13}
\tightlist
\item
  设I是一个预有序集，并设 \((\mathbf{E}_{\alpha})_{\alpha \in \mathbf{I}}\) 是由I指标化的一族集合。对于每一对 \((\alpha, \beta)\) （I 的指标，满足 \(\alpha \leqslant \beta\) ），让 \(f_{\alpha \beta}\) 是从 \(\mathbf{E}_{\beta}\) 到 \(\mathbf{E}_{\alpha}\) 的映射，并假设关系 \(\alpha \leqslant \beta \leqslant \gamma\) 蕴含 \(f_{\alpha \gamma} = f_{\alpha \beta} \circ f_{\beta \gamma}\) .
\end{enumerate}

设 G 为集合族 \((\mathrm{E}_{\alpha})_{\alpha\in\mathbf{I}}\) 的乘积。设 E 为 G 中满足所有关系 \(pr_{\alpha}x = f_{\alpha\beta}(pr_{\beta}x)\) 的所有元素 x 构成的子集，其中指标对 \(\alpha, \beta\) 使得 \(\alpha \leqslant \beta\) 取遍所有指标对。集合 E 称为族 \((\mathrm{E}_{\alpha})_{\alpha\in\mathbf{I}}\) 关于映射族 \((f_{\alpha\beta})\) 的逆向极限，且限制 \(f_{\alpha}\) （即 \(pr_{\alpha}\) 在 E 上的限制）称为 E 到 \(E_{\alpha}\) 的典范映射。我们有 \(f_{\alpha} = f_{\alpha\beta} \circ f_{\beta}\) 每当 \(\alpha \leqslant \beta\) 。我们记 \(E = \varprojlim(E_{\alpha}, f_{\alpha\beta})\) ，或简称为 \(E = \varprojlim E_{\alpha}\) 当无混淆风险时。滥用语言地，称对 \(((E_{\alpha}), (f_{\alpha\beta}))\) 为相对于I的集合的逆向系统。

应当指出，即使所有 \(E_{\alpha}\) 都非空且所有映射 \(f_{\alpha\beta}\) 是满射的，E 也可以为空集。

对于每一个 \(\alpha \in I\) 让 \(g_{\alpha}\) 是从集合 \(E'\) 到 \(E_{\alpha}\) 的映射，使得关系 \(\alpha \leqslant \beta\) 蕴含 \(f_{\alpha \beta} \circ g_{\beta} = g_{\alpha}\) 。则存在唯一的映射 \(g\) （从 \(E'\) 到 \(E\) ）使得 \(g_{\alpha} = f_{\alpha} \circ g\) 对于所有 \(\alpha \in I\) . 要使 \(g\) 为单射，必要且充分条件是：对于每一对不同的元素 \(x', y'\) （属于 \(E'\) ），应当存在 \(\alpha \in I\) 使得 \(g_{\alpha}(x') \neq g_{\alpha}(y')\) .

让 \((\mathbf{A}_{\alpha},\varphi_{\alpha \beta})\) 并且 \((\mathbf{B}_{\alpha},\psi_{\alpha \beta})\) 是两个相对于同一指标集 I 的集合的逆向系统。设 \(\mathbf{A} = \varprojlim (\mathbf{A}_{\alpha},\varphi_{\alpha \beta}),\mathbf{B} = \varprojlim (\mathbf{B}_{\alpha},\psi_{\alpha \beta})\) ，并且对于每个 \(\alpha \in I\) 让 \(\varphi_{\alpha}\) （相应地 \(\psi_{\alpha}\) ）为 A 到 \(\mathbf{A}_{\alpha}\) （相应地，B 到 \(\mathbf{B}_{\alpha}\) ）的典范映射。对于每一个 \(\alpha \in I\) 让 \(u_{\alpha}\) 是从 \(\mathbf{A}_{\alpha}\) 到 \(\mathbf{B}_{\alpha}\) 的映射，使得 \(\psi_{\alpha \beta}\circ u_{\beta} = u_{\alpha}\circ \varphi_{\alpha \beta}\) 每当 \(\alpha \leqslant \beta\) 。族 \((u_{\alpha})\) 被称为 \((\mathbf{A}_{\alpha},\varphi_{\alpha \beta})\) 到 \((\mathbf{B}_{\alpha},\psi_{\alpha \beta})\) 的映射逆向系统。在这些条件下，存在唯一的映射 \(u:\mathbf{A}\to \mathbf{B}\) 使得 \(\psi_{\alpha}\circ u = u_{\alpha}\circ \varphi_{\alpha}\) 对于所有 \(\alpha \in I\) . 映射 \(u\) 称为 \(u_{\alpha}\) 的逆向极限，记作 \(u = \varprojlim u_{\alpha}\) 当没有混淆风险时。设 \((\mathbf{C}_{\alpha}, \theta_{\alpha\beta})\) 为相对于 I 的另一个集合逆向系统，设 \(v_{\alpha}\) 是 \((\mathbf{B}_{\alpha}, \psi_{\alpha\beta})\) 到 \((\mathbf{C}_{\alpha}, \theta_{\alpha\beta})\) 的映射逆向系，并且让 \(v = \varprojlim v_{\alpha}\) 。于是我们有 \(\lim (v_{\alpha} \circ u_{\alpha}) = v \circ u\) .

设 J 是 I 的一个共尾子集，并假设 J 是有向的（关于由 I 诱导的序）。如果 \(((E_{\alpha})_{\alpha\in I}, (f_{\alpha\beta})_{\alpha\in I,\beta\in I})\) 是相对于 I 的一个集合逆向系，其逆向极限为 E，则 \(((E_{\alpha})_{\alpha\in J}, (f_{\alpha\beta})_{\alpha\in J,\beta\in J})\) 也是相对于 J 的一个集合逆向系。设 \(E'\) 为其逆向极限，并令 \(f_{\alpha}'\) 是从 \(E'\) 到 \(E_{\alpha}\) 的典范映射，其中 \(\alpha\in J\) 。对于每一个 \(x\in E\) 让 \(g(x)=(f_{\alpha}(x))_{\alpha\in J}\in E'\) （其中 \(f_{\alpha}\) 表示从E到 \(E_{\alpha}\) 的典范映射）。则 g 是 E 到 \(E'\) 的双射，借助它 \(E'\) 通常与 E 等同。

